\lecture{11}{29. September 2025}{Internal Flows: Flows in Pipes and Ducts}

\begin{problem}{8.20}
  For a turbulent flow of a fluid in \qty{0,6}{\m} diameter pipe, the velocity \qty{0,15}{\m} from the wall is \qty{2,7}{\m\per\s}.
  Estimate the wall shear stress using the 1/7\textsuperscript{th} expression for the velocity profile.
\end{problem}

\begin{problem}{8.24}
  A pipe runs for an elevation of \qty{45}{\m} to an elevation of \qty{115}{\m}.
  The inlet pressure is \qty{8,5}{\MPa} and the head loss is \qty{6,9}{\kJ\per\kg}.
  Calculate the outlet pressure for (a) the inlet at the \qty{45}{\m} elevation and (b) the inlet at the \qty{115}{\m} elevation.
\end{problem}

\begin{problem}{8.28}
  Water flows in a smooth \qty{0,2}{\m} diameter pipeline that is \qty{65}{\m} long.
  The Reynolds number is \num{e6}.
  Determine the flow rate and pressure drop.
  After many years of use, minerals deposited on the pipe cause the same pressure drop to produce only one-half the original flow rate.
  Estimate the size of the relative roughness elements for the pipe.
\end{problem}
